\section{Introduction}
\label{bl:sec:introduction}

A nonzero section of an adjoint divisor gives an effective representative;
abundance requires enough sections to define a morphism. Adjunction
suggests looking first on the support of the zero divisor, where a
multiple may already be generated. We prove that, when this support lies
in the coefficient-one boundary, generation on the entire reduced support
forces positive ambient Iitaka dimension.

\begin{theorem}[Supported boundary lifting]\label{bl:thm:supported}
Let $(V,C)$ be a normal projective $\mathbb Q$-factorial dlt pair over an
algebraically closed field of characteristic zero, with $C$ an effective
rational boundary, and put $A=K_V+C$. Suppose that
\[
  0\ne G\geq0,\qquad G\sim qA,\qquad
  \operatorname{Supp}G\subseteq\operatorname{Supp}\lfloor C\rfloor,
\]
where $G$ is Cartier and $q>0$ is a sufficiently divisible integer.
If $\mathcal O_V(G)|_{G_{\mathrm{red}}}$ is semiample, then
$\kappa(V,A)>0$.
\end{theorem}

The restriction is a line bundle on the whole reduced scheme
$S=G_{\mathrm{red}}$. Sections on its separate components must agree
along their intersections. The support inclusion may be strict: in
$C=S+H+C_0$, the extra coefficient-one divisor $H$ may meet $S$ and
any of its strata, while $C_0$ has coefficients less than one.
No additional nefness hypothesis is needed. Curves in $S$ have
nonnegative $G$-degree by semiampleness, and curves outside $S$ have
nonnegative degree by effectivity. The rational coefficients of $C$
may be arbitrary in $[0,1]$.

The first application separates abundance from the task of obtaining a
first section.
\begin{theorem}[Abundance after nonvanishing]
\label{thm:abundance-after-nonvanishing}
Let $(X,\Delta)$ be a projective log canonical pair of dimension at most
four over $\C$, with
effective rational boundary. If $D=K_X+\Delta$ is $\mathbb Q$-Cartier
and nef and $\kappa(X,D)\geq0$, then $D$ is semiample. If
$\kappa(X,D)=0$, then $D\sim_{\mathbb Q}0$.
\end{theorem}

The proof uses lower-dimensional abundance and ordinary minimal models
for effective log canonical pairs. Theorem~\ref{thm:abundance-after-nonvanishing}
is independent of the existence of a first section in dimension four.
It applies to arbitrary rational boundary coefficients and supplies a
qualitative generated Cartier multiple, with no uniform index.
Section~\ref{bl:sec:fields} transfers the same conclusion to every
algebraically closed field of characteristic zero.

The same supported argument gives a reduction in all dimensions.
\begin{theorem}[Reduction to smooth canonical nonvanishing]
\label{bl:thm:conditional}
Assume that every smooth connected projective complex variety $W$ with
pseudo-effective $K_W$, in every dimension, has a nonzero section of
$mK_W$ for some $m>0$. Let $(X,\Delta)$ be a normal projective
log canonical pair over an algebraically closed field of characteristic
zero, with $\Delta$ an effective rational boundary and
$D=K_X+\Delta$ a nef $\Q$-Cartier divisor. Then $D$ is semiample.
\end{theorem}
The premise concerns pseudo-effective canonical divisors and includes
those that are not nef. Hashizume's reduction supplies nonvanishing
and ordinary log minimal models under this premise; it does not
supply the semiampleness proved here.

A separate companion provides the first section needed in dimension four.
\begin{corollary}[Fourfold log abundance]\label{thm:full-abundance}
Let $(X,\Delta)$ be a normal projective log canonical pair of dimension
at most four over an algebraically closed field of characteristic zero.
If $\Delta$ is an effective rational boundary and $D=K_X+\Delta$ is
$\Q$-Cartier and nef, then $D$ is semiample. If $\kappa(X,D)=0$,
then $D\sim_{\Q}0$.
\end{corollary}
This resolves rational log abundance positively in the stated range,
using Theorem~\ref{thm:abundance-after-nonvanishing} and the nonvanishing
result in OpenAI, \emph{Fourfold nonvanishing by minimal metrics and moving jets}
(September 27, 2026), Corollary~1.2
\cite{OpenAINonvanishing2026}. The proof of the supported theorem and
of abundance after nonvanishing uses no result from that companion.

\subsection*{Boundary deformation and extension}

The relation between boundary geometry and ambient sections already
appears in dimension three. Kawamata's alternative proof of Miyaoka's
numerical-dimension-one theorem constructs compatible infinitesimal
boundary deformations using finite covers and Hodge theory, and then
produces a moving family \cite[Section 4]{Kawamata1992}. The question
of how infinitesimal boundary data yield ambient positivity also
motivates the lifting problem considered here.

Log abundance through dimension three rests on the surface and threefold
theorems, including the corrected Keel--Matsuki--McKernan argument
\cite{KMM1994,KMM2004,Fujino2012} and semi-log-canonical gluing
\cite{Fujino2000}. Liu--Xu proved abundance after nonvanishing through
dimension five under the additional numerical-dimension bound
$\nu\leq1$, together with a corresponding conditional reduction from
smooth canonical nonvanishing in dimension $d-1$ and in numerical
dimension one in dimension $d$
\cite[Theorems~5.1 and~5.4]{LiuXu2025}. Theorem~\ref{thm:abundance-after-nonvanishing}
removes that numerical-dimension bound in dimension at most four;
the lifting theorem itself holds in every dimension.

On the analytic side, Demailly--Hacon--P\u{a}un proved extension for
plt pairs without requiring an ample boundary part and explained its
connection with nonvanishing and good minimal models
\cite[Theorem 1.7 and Corollary 1.8]{DHP2013}. Their nef plt
corollary gives restriction surjectivity in the presence of a suitable
effective supported representative. Chan--Choi subsequently removed
one support-containment restriction in the log-smooth plt setting
\cite[Theorem 1.6.1]{ChanChoi2021}. Theorem~\ref{bl:thm:supported}
starts with semiampleness on an entire reducible support and allows
additional coefficient-one components to meet it. Compatibility at
those intersections, and obstructions concentrated on smaller strata,
are central to its proof.

The passage to abundance also needs distinct results about gluing and
minimal models. Birkar's effective log minimal model results provide
the birational models used in the induction
\cite{Birkar2010,Birkar2011}. Fujino--Gongyo show that semiampleness
on the normalization of a semi-log-canonical pair, with its conductor
boundary, descends to the original pair
\cite[Theorem 4.3]{FujinoGongyo2014}. Their nef log-abundance theorem
requires abundance on the ambient variety and on the normalizations of
all log canonical centers \cite[Theorem 4.2]{FujinoGongyo2014}.
Both requirements enter the proof of Theorem~\ref{thm:abundance-after-nonvanishing}.

\subsection*{The obstruction on a finite neighborhood}

After replacing $G$ by a multiple, its generated restriction gives
$f:S\to\mathbb P^\ell$ with
$\mathcal O_V(G)|_S\simeq f^*\mathcal O(1)$. We want to extend
functions and sections through successive infinitesimal neighborhoods
of $S$. There is an immediate difficulty: $f$ itself need not extend
to any neighborhood. The proof must therefore measure the lifting
obstruction using only the morphism on the reduced boundary.

Pass to the bundle of nonzero frames of $\mathcal O_V(G)$, where the
pulled-back line bundle has a tautological trivialization. Scalar
multiplication of a frame gives an action of $\mathbb G_m$. A function
or form has weight $a$ if this action multiplies it by the character
$\lambda^a$. Two cyclic-cover constructions, retaining all components
of each normalized cover, produce a reduced Cartier divisor $E=(y=0)$.
After an equivariant resolution there is a logarithmic top form $\sigma$
whose residue along the reduced resolved boundary has positive weight. Poles over the extra
boundary $H$ remain present, including above its intersections with $S$.

Locally, the finite neighborhoods are defined by $(y^k)$.
At the first stage where a lift from $y^k=0$ to $y^{k+1}=0$ is not yet
known, differences of local lifts are multiples of $y^k$. Their
coefficients form a cohomology class on $E$: this is the obstruction.
Because products of two such errors vanish at that order, the connecting
map on functions satisfies the Leibniz rule. Thus the obstruction is a
derivation. A suitable character part of the residue map embeds its
cohomology group into the cohomology of logarithmic residues.

The identity $\mathcal O_V(G)|_S=f^*\mathcal O(1)$ lifts $f$ to a
map between the bundles of nonzero frames. On a chart trivializing
$\mathcal O(1)$, the target of this lifted map has projective-base
coordinates and an invertible frame scalar. Changing the base coordinates
gives a horizontal direction; rescaling that scalar gives the vertical
direction. The proof first eliminates the
horizontal part of the obstruction, then uses the nonzero weight of the
frame action to eliminate the vertical part. This distinction explains
why positivity of the residue weight is useful.

To retain classes supported on singular fibers or smaller strata, the
residue comparison is made in a filtered direct image of a Hodge module
on the graph of the boundary morphism. Its lowest filtration term is the
whole coherent residue cohomology sheaf. Its first symbol records the
effect of differentiating in coordinates on the target of the lifted map.
Saito's strictness and
Kodaira--Saito vanishing provide the required control of these filtered
objects \cite{Saito1988,Saito1990,Popa2016}; restricting to the smooth
locus of a variation of Hodge structure would discard some of the
obstructions we need to test.

A local calculation on the graph expresses the derivation and this
first symbol in a single residue identity. The symbol has positive frame
weight. If its horizontal part were nonzero, a finite cover transverse
to that part, together with the adjugate of its differential, would
produce a nonzero map from an ample line bundle into the kernel of the
first symbol. Kodaira--Saito vanishing excludes that map. The remaining
vertical term is the actual Euler derivative of frame scaling. The
residue identity now says that a nonzero scalar, determined by its weight,
times the vertical obstruction is zero. Hence that obstruction vanishes
as well. Applying the same identity
to arbitrary boundary functions removes the complete obstruction. This
proves lifting at every finite order.

Taking finite-cover invariants and the degree-zero part of the frame
action descends the lifted layers to divisorial sheaves on $V$.
These layers recur with positive twists of $\mathcal O_{\mathbb P^\ell}(1)$.
Their accumulated sections grow while higher cohomology stays bounded;
the loss in passing from a finite neighborhood to ambient sections is
bounded by a fixed cohomology group. Hence the section spaces of
multiples of $G$ are unbounded, proving positive Iitaka dimension.
If the boundary morphism has zero-dimensional image, the growing number
of layers supplies this increase. The argument uses finite filtrations
and does not require convergence of a formal neighborhood.

\subsection*{From lifting to abundance and other boundary methods}

For Theorem~\ref{thm:abundance-after-nonvanishing}, lower-dimensional abundance and
conductor gluing generate the adjoint on the whole reduced floor.
Theorem~\ref{bl:thm:supported} then excludes a nonzero effective
representative in the case $\kappa=0$. For positive Iitaka dimension,
effective minimal models handle the generally non-nef adjoints on the
Iitaka fibers. Numerical triviality there gives abundance, including
on log canonical centers, and the preceding semiampleness and descent
results apply.

The complex lifting argument also works under a weaker formulation:
the canonical and boundary divisors are individually $\mathbb Q$-Cartier,
the pair is log canonical, and it is klt off the reduced floor. This
formulation permits transfer to other algebraically closed fields of
characteristic zero. One descends finitely many defining data to a countable algebraically closed
field embeddable in $\mathbb C$ and transfers sections and generation
by base change of the original line bundle and its evaluation map.
No preservation of $\mathbb Q$-factoriality under that descent is needed.

A further method illuminates the boundary obstruction with different
hypotheses. A complete conormal--period argument proves that a smooth
projective complex fourfold with $K_X$ nef and $\kappa(K_X)=0$ cannot
have $\nu(K_X)=2$, without an Euler-characteristic hypothesis.
Its period equation and trace tests at exceptional parameters kill
successive obstructions; a quadratic count of invariant layers then
forces growth. This method is not needed for
Theorem~\ref{bl:thm:supported}.

\subsection*{Reading the proof}
Section~\ref{bl:sec:covers} constructs the two covers and a split residue
insertion. Section~\ref{bl:sec:hodge} identifies the lowest filtered
piece and proves the symbol-kernel vanishing. Section~\ref{bl:sec:frames}
separates horizontal symbols from the actual Euler action, and
Section~\ref{bl:sec:lifting} uses these two conclusions to lift and
count all finite layers. Section~\ref{bl:sec:abundance} proves the
abundance implications over $\C$; Section~\ref{bl:sec:fields} returns
the resulting sections and generated multiples to the original field.
The independent curve-base method is developed in
Section~\ref{nu2:sec:method}.

\subsection*{Conventions and the complex formulation}
All boundaries are effective with rational coefficients in $[0,1]$.
An adjoint is an actual $\Q$-Cartier divisor, and $\sim_{\Q}$ denotes
rational linear equivalence. A rational Cartier divisor is semiample
if a positive Cartier multiple is generated by global sections; its
Iitaka dimension is the dimension of the images of all sufficiently
divisible nonempty complete systems, with value $-\infty$ if no such
system exists. For a nef divisor $D$ on a projective $d$-fold,
\[
 \nu(D)=\max\{j:D^jH^{d-j}>0\}
\]
for an ample divisor $H$; these intersections may be computed on a
smooth resolution. We use log discrepancies, so coefficient-one
divisors have log discrepancy zero. Standard conventions for lc,
klt, dlt and dlt adjunction follow \cite{Kollar2013}.

The complex proof uses only individual Cartier witnesses, rather than
preservation of $\Q$-factoriality under a field extension.
\begin{proposition}[Complex supported assertion]
\label{bl:prop:complex-supported}
Let $V$ be normal and projective over $\C$, and let $C$ be an
effective rational boundary. Suppose that $K_V$ and every component
of $C$ are individually $\Q$-Cartier, that $(V,C)$ is lc, and that
it is klt outside $\Supp\lfloor C\rfloor$. If a nonzero effective
Cartier divisor $G\sim q(K_V+C)$ has support in
$\Supp\lfloor C\rfloor$, for a sufficiently divisible positive
integer $q$, and $\OO_V(G)|_{G_{\mathrm{red}}}$ is semiample on the
whole reduced scheme, then $\kappa(V,K_V+C)>0$.
\end{proposition}
This proposition is proved in Sections~\ref{bl:sec:covers}--\ref{bl:sec:lifting}.
Its hypotheses follow from those of Theorem~\ref{bl:thm:supported}
over $\C$, and its individual Cartier requirements are the finite
witnesses used in the final field transfer.
